% HD-F1-T31. GPT-6 Astra Pro; /root/transport; 22 November 2026 UTC.
% Effort inherited; exact setting not exposed.
\section[HD-F1-T31: at most five symplectic quotient branches]{HD-F1-T31: punctured actual factor tuples give at most five quotient branches and multiplicity at least fifteen}
\label{sec:f1-symplectic-punctured-tuples}

Retain the specified scalar-centre symplectic-type normal subgroup
\(R\lhd G\) of T17 and the original remaining T2 source.
Use its actual evaluation isomorphism
\[
 W\otimes M\xrightarrow{\sim}V_{\mathbf C},\qquad
 d=\dim W=2^a\ (a\ge1),\quad k=\dim M,\quad m=dk\ge27.
 \tag{F1T31.1}
\]
The first nine multiplicities have already been excluded, so
\(k\ge10\). Put \(C_R=C_G(R)\), let \(n\) be the number
of original inertia entries outside \(C_R\), and set
\[
 q=d\left\lceil\frac5d\right\rceil.
 \tag{F1T31.2}
\]
Then the following are necessary for every remaining source:
\[
 \begin{gathered}
                    3\le n\le5,\qquad k\ge15,\\
 \begin{array}{c|c}
 n&\text{integral original-source inequality}\\ \hline
 3&m+12\ge2(b-3)q\\
 4&m+24\ge2(b-4)q\\
 5&m+48\ge3(b-5)q.
 \end{array}
 \end{gathered}
 \tag{F1T31.3}
\]
The actual quotient cover by \(C_R\) has exactly \(n\)
branch points, and hence at most two branch parameters after
projective changes of the base. These are necessary restrictions
on the original source; no existence assertion for a remaining
integer or tuple is made.

\paragraph{1. Puncturing an irreducible product-scalar tuple.}
Let \(B_1,\ldots,B_s\) act irreducibly on a complex space
\(U\) of dimension \(h>1\), and suppose their ordered
product is \(z1_U\). For every index \(i\), the actual map
\[
 \begin{split}
 F_i:U&\longrightarrow\bigoplus_{j\ne i}\operatorname{Im}(B_j-1),\\
 w&\longmapsto((B_j-1)w)_{j\ne i}
 \end{split}
 \tag{F1T31.4}
\]
is injective. Indeed its kernel is fixed by all \(B_j\)
with \(j\ne i\). For a vector \(w\) in that kernel,
the product identity gives
\[
 B_iw=z(B_1\cdots B_{i-1})^{-1}w=zw.
\]
Thus the kernel is invariant under every factor. It cannot
be all of \(U\), since then every factor is scalar and
\(h>1\) contradicts irreducibility. It is zero, proving
\[
 \sum_{j\ne i}\operatorname{rank}(B_j-1)\ge h,
 \qquad
 \sum_j\operatorname{rank}(B_j-1)
          \ge h+\max_j\operatorname{rank}(B_j-1).
 \tag{F1T31.5}
\]
This proof retains the scalar \(z\); it does not replace
the product relation by product one.

The T20/T30 intertwiners give the original identities
\(g_j=A_j\otimes T_j\). Their two tuples are irreducible:
an invariant subspace of \(M\) gives the invariant subspace
\(W\otimes M'\) of the original representation; a word in
the \(g_j\) equal to \(r\in R\) gives a word in the
\(A_j\) equal to a scalar times \(\rho_W(r)\).
The latter forces every invariant subspace in \(W\) to
be \(R\)-invariant. The original ordered product one gives
scalar ordered products on both actual factors.

For entries inside \(C_R\), take \(A_j=1\).
For each outside entry, first normalize \(A_j\) at a
largest eigenspace and change \(T_j\) by the inverse
scalar, as in T30. Next choose a largest eigenvalue
\(\mu_j\) of this actual \(T_j\), and put
\(\widehat T_j=\mu_j^{-1}T_j\). Thus the retained identity is
\[
                 g_j=\mu_j A_j\otimes\widehat T_j.
 \tag{F1T31.6}
\]
Inside \(C_R\) the same identity holds with \(A_j=1\).
Both independently normalized tuples still have scalar
ordered products and the same invariant subspaces. Their
factors are semisimple by the finite-projective-order
argument of T20. Define
\[
 x_j=\frac{\operatorname{rank}(A_j-1)}d,
 \qquad
 y_j=\frac{\operatorname{rank}(\widehat T_j-1)}k.
 \tag{F1T31.7}
\]
These are minimum eigenspace codimensions divided by the
factor dimensions. The actual symplectic intertwiner bound
of T17/T20 gives \(x_j\ge1/4\) outside \(C_R\), and
\(x_j=0\) inside. Applying (F1T31.5) gives
\[
 \sum_{j\ne i}x_j\ge1,
 \qquad \sum_{j\ne i}y_j\ge1
                  \quad\text{for every }i.
 \tag{F1T31.8}
\]
No separate factor is declared to be a curve coefficient
system or a replacement for the original sixfold corner.

\paragraph{2. A tensor bound for every original scalar phase.}
For semisimple matrices with normalized minimum eigenspace
codimensions \(x,y\), every scalar multiple of their
tensor product has moving-space fraction at least
\[
       F(x,y)=\max\{x,y,x+y-2xy\}.
 \tag{F1T31.9}
\]
To prove this, arrange their normalized eigenspace
multiplicities in decreasing order, \((p_i)\) and \((t_i)\).
An eigenspace for a specified tensor eigenvalue pairs
distinct eigenvalues of the two factors, and its dimension
is the corresponding partial matching sum. Adding zero
terms and exchanging two out-of-order pairs bounds that
sum by \(\sum_i p_it_i\). This is at most
\[
 p_1t_1+\Bigl(\sum_{i>1}p_i\Bigr)
                   \Bigl(\sum_{i>1}t_i\Bigr)
       =(1-x)(1-y)+xy.
\]
It is also at most each of \(p_1,t_1\). Subtracting
these three bounds from one proves (F1T31.9), for the
literal eigenvalue one in (F1T31.6) and every \(\mu_j\).

The function \(F\) is nondecreasing in each variable.
When both variables are at most \(1/2\), it equals
\(x+y-2xy\); when at least one is at least \(1/2\),
it equals \(\max(x,y)\). In particular, for \(x\ge1/4\),
\[
                  F(x,y)\ge\frac14+\frac12y.
 \tag{F1T31.10}
\]
For \(y\le1/2\) this follows by decreasing \(x\) to
\(1/4\). For \(y\ge1/2\), use
\(F(x,y)\ge y\ge1/4+y/2\).

Write \(r_j=\operatorname{rank}(g_j-1)\) and set
\[
 H=\sum_{g_j\in C_R}r_j,\qquad W=H/m.
 \tag{F1T31.11}
\]
Each nonidentity inside entry has rank a positive multiple
of \(d\), and the retained original gate makes its rank
at least five. Therefore
\[
                  H\ge(b-n)q.
 \tag{F1T31.12}
\]
The minimum eigenspace codimension of an inside \(T_j\)
is no larger than its actual codimension at eigenvalue one,
so \(\sum_{\rm inside}y_j\le W\). Put
\(Y=\sum_{\rm outside}y_j\) and
\(t=\max_{\rm outside}y_j\). Puncturing at an outside
index attaining \(t\) gives
\[
        Y-t+W\ge1,
        \qquad Y\ge\frac{n}{n-1}(1-W).
 \tag{F1T31.13}
\]
The latter inequality follows from \(t\ge Y/n\).
It is valid also for \(W\ge1\); its right side is
then nonpositive, and no equality or fixed value for
the inside ranks has been imposed.

Summing (F1T31.10) and retaining the exact inside rank gives
\[
 \begin{split}
 \sum_j r_j
  &\ge m\left(\frac n4+\frac12Y+W\right)\\
  &\ge\frac{n(n+1)}{4(n-1)}m
             +\frac{n-2}{2(n-1)}H.
 \end{split}
 \tag{F1T31.14}
\]

\paragraph{3. At most five outside-centralizer inertia entries.}
The actual quotient \(G/C_R\) contains \(R/Z_R\),
which is noncyclic. Its original generating product-one
tuple therefore has at least three nonidentity entries,
as proved in T17. Thus \(n\ge3\).
If \(n\ge9\), the outside rank bound alone gives
\(2m+6\ge9m/4\), contrary to \(m\ge27\).
For \(n=6,7,8\), the coefficient of \(m\) in
(F1T31.14) is respectively
\[
                   \frac{21}{10},\quad\frac73,
                         \quad\frac{18}{7},
\]
each larger than two. Using \(b\ge12\) and \(q\ge5\),
the remaining term in that bound is respectively at least
\[
                    12,\quad\frac{125}{12},
                              \quad\frac{60}{7},
\]
each larger than six. This contradicts the original
identity \(\sum r_j=2m+6\). Hence \(3\le n\le5\).
For \(n=5\), (F1T31.14) is
\[
                \sum_jr_j\ge\frac{15}{8}m+\frac38H.
 \tag{F1T31.15}
\]
Together with (F1T31.12), this proves the third line
of (F1T31.3) without restricting any outside spectrum.

\paragraph{4. Three outside entries.}
The three punctured \(A\)-inequalities give
\(2\sum_{\rm outside}x_j\ge3\).
The original tensor rank is at least \(mx_j\), so
\[
                  \sum_jr_j\ge\frac32m+H.
 \tag{F1T31.16}
\]
Equating the left side with \(2m+6\) and using
(F1T31.12) proves the first line of (F1T31.3).

\paragraph{5. Four outside entries: the entire joint optimization.}
We prove the uniform inequality
\[
                  \sum_jr_j\ge\frac74m+\frac12H.
 \tag{F1T31.17}
\]
First replace each of the four outside \(x_j\) by
\(\min(x_j,1/2)\). This decreases \(F\) and preserves
every punctured \(A\)-inequality. Indeed a sum of the
other three terms is unchanged if none was truncated;
if one was truncated, the new sum is at least
\(1/2+1/4+1/4=1\).

For these truncated variables put \(a_j=4x_j-1\).
They lie in \([0,1]\) and satisfy
\(\sum_{j\ne i}a_j\ge1\). Reorder them as
\(a_1\ge a_2\ge a_3\ge a_4\), and put
\(S=a_2+a_3+a_4\ge1\). The vector
\[
       (t,t,s,r)=\frac1S(a_2,a_2,a_3,a_4)
\]
is componentwise at most the original vector. Here
\(t\ge s\ge r\ge0\), \(t+s+r=1\), and it has
the exact convex decomposition
\[
 \begin{split}
 (t,t,s,r)={}&(t-s)(1,1,0,0)\\
 &+2(s-r)(1/2,1/2,1/2,0)\\
 &+3r(1/3,1/3,1/3,1/3).
 \end{split}
 \tag{F1T31.18}
\]
The coefficients are nonnegative and sum to one.
For fixed \(y\), the function \(F(x,y)\) is affine
in \(x\in[1/4,1/2]\): it is
\(y+(1-2y)x\) for \(y\le1/2\), and the constant
\(y\) for \(y\ge1/2\). Consequently it is enough
to check the following three outside \(x\)-vectors,
with the corresponding permutation of the actual \(y_j\):
\[
 (1/2,1/2,1/4,1/4),\quad
 (3/8,3/8,3/8,1/4),\quad
 (1/3,1/3,1/3,1/3).
 \tag{F1T31.19}
\]
This is an exact convex reduction of an inequality on
the actual factor data, not a replacement of any matrices.

For the first vector put
\(p=(y_1-1/2)_++(y_2-1/2)_+\) and \(z=y_3+y_4\).
Its outside cost is at least \(3/2+p+z/2\).
On the other hand (F1T31.13) implies
\[
 1-W\le\sum_jy_j-\max_jy_j
       \le\min(y_1,y_2)+z\le1/2+p/2+z.
\]
It follows that \(p+z/2\ge1/4-W/2\).
After adding the exact inside cost \(W\), the total
normalized cost is at least \(7/4+W/2\).

For the second vector put \(P=y_1+y_2+y_3\), \(z=y_4\).
Its outside cost is at least \(11/8+P/4+z/2\).
Since \(\max_jy_j\ge P/3\), the punctured inequality gives
\((2/3)P+z\ge1-W\). Therefore the total cost is at least
\[
 \frac{11}{8}+\frac38(1-W)+W
             =\frac74+\frac58W
             \ge\frac74+\frac12W.
\]
For the third vector, its outside cost is at least
\(4/3+(\sum_jy_j)/3\). Puncturing gives
\(\sum_jy_j\ge(4/3)(1-W)\), so the total is at least
\[
             \frac{16}{9}+\frac59W
                 \ge\frac74+\frac12W.
\]
All these inequalities hold for every \(W\ge0\).
This proves (F1T31.17) on the full polytope and hence
on the actual eigenspace data. It gives the middle
line of (F1T31.3).

\paragraph{6. Explicit integral consequences.}
Using \(b\ge12\) in (F1T31.3), with the exact
integer \(q\), gives the following lower bounds for \(k\):
\[
 \begin{array}{c|ccc}
 d&n=3&n=4&n=5\\ \hline
 2&48&36&39\\
 4&33&26&30\\
 8&17&13&15\\
 16&18&15&18\\
 32&18&16&20\\
 2^a\ge64&18&16&21.
 \end{array}
 \tag{F1T31.20}
\]
For example the three bounds are respectively the
ceilings of \((18q-12)/d\), \((16q-24)/d\), and
\((21q-48)/d\). The last row follows for every power
of two at least sixty-four, not from a finite sample.
These bounds first give \(k\ge13\), with \(k=13\)
or \(14\) confined to \(d=8,n=4\). The next paragraph
excludes those two values.

\paragraph{7. The last two multiplicities.}
Suppose \(d=8,n=4\) and \(k\) is thirteen or fourteen.
The exact symplectic eigenspace gap of T30 says that
the outside minimum codimensions \(u_j=dx_j\) are
two or at least three. They are integers. Clip them
at four, preserving every punctured sum of the other
three and decreasing \(F\), by the proof of paragraph5.
The resulting integers are in \(\{2,3,4\}\), and their
three smallest members have sum at least eight.
If the two smallest are both two, the other two are
both four. Otherwise the vector dominates a permutation
of \((2,3,3,3)\). Thus every case dominates one of
\[
                         (2,2,4,4),\qquad(2,3,3,3).
 \tag{F1T31.21}
\]
The first vector already gives an outside cost at least
\(12k\). Since \(H\ge8q=64\), its total is at least
\(12k+64>16k+6\) for both proposed values of \(k\).

For the second vector let
\(S=\sum_{\rm outside}\operatorname{rank}(\widehat T_j-1)\).
The tensor bound (F1T31.9) gives an outside cost at least
\(11k+2S\): the coefficient \(d-2u_j\) is at least
two in each term. The punctured \(T\)-inequality gives
\[
                  S\ge\frac43\left(k-\frac H8\right).
\]
This remains valid when \(H/8>k\). Consequently the
total original cost is at least
\[
                  \frac{41k+2H}{3}
                     \ge\frac{41k+128}{3}>16k+6
                   \quad(k=13,14).
 \tag{F1T31.22}
\]
The two differences from \(16k+6\) are respectively
\(58-4k\) and \((110-7k)/3\), both positive at
the two stated integers. This excludes the last two
multiplicities and proves \(k\ge15\).

\paragraph{8. Actual quotient cover and original corner maps.}
Let \(X\to X/G=\mathbf P^1\) be the original full
connected cover used for the faithful coefficient group.
Since \(C_R\) is normal, the actual finite quotient gives
\[
 X\xrightarrow{\alpha}X/C_R
       \xrightarrow{\beta}\mathbf P^1,
 \qquad \deg\alpha=|C_R|,
 \quad\deg\beta=|G/C_R|.
 \tag{F1T31.23}
\]
The inertia of \(\beta\) is the image of the original
inertia in \(G/C_R\). It is nontrivial exactly at the
\(n\) original outside-centralizer branch values.
The map \(G\to\operatorname{PGL}(W)\) furnished by the
actual intertwiners has kernel exactly \(C_R\): a scalar
intertwiner centralizes every \(R\)-matrix, and conversely
Schur's lemma makes every such intertwiner scalar.
Thus this quotient is precisely the projective action
on the specified symplectic factor.

On a fixed finite-group Hurwitz component, passing to
\(X/C_R\) and forgetting the unramified branch labels
is an actual morphism of cover families. Fixing its
\(n\) branch values leaves only finitely many possible
covering monodromies: each is a homomorphism from the
finitely generated punctured-sphere group to the fixed
finite group \(G/C_R\). Removing projective changes
of the base therefore bounds this quotient-cover family
by \(n-3\le2\) parameters. This statement does not
bound the full original period image by two.

The quotient pullback and norm satisfy
\(\alpha_*\alpha^*=[|C_R|]\) on the quotient Jacobian,
and the cup product obeys
\(\langle\alpha^*u,\alpha^*v\rangle_X
 =|C_R|\langle u,v\rangle_{X/C_R}\).
These are the push--pull and Poincare-biextension maps
for the displayed quotient, not an identification with
the original six-dimensional corner. That corner,
its rational group-algebra projector, restricted theta
polarization, coefficient local system and variation
are retained verbatim. The evaluation (F1T31.1),
the scalar identities (F1T31.6), and extension from
\(K\) to \(\mathbf C\) preserve the original matrix
ranks. On its Chow/variational, Poincare, de Rham and
Adams realizations the comparison is the identity with
polarization multiplier one, as in F01--F04. The proof
uses no classification of other normal subgroups and
constructs no additional detector component.
